\section{Máximo de funciones}

\subsection{Li-Chao Tree}
La construcción inicial toma $\mathcal{O}(N \log N)$. Dado un conjunto $A \subset \mathbb{R}$ con $|A|=N$ y $M$ funciones que cumplen la \textit{transcending property}, la estructura permite agregar una nueva función en $\mathcal{O}(\log N)$ por operación. Además, para cualquier $a \in A$, permite responder $Q$ consultas de la forma $\max_i f_i(a)$ en $\mathcal{O}(\log N)$ cada una. Complejidad total: $\mathcal{O}((N + M + Q) \log N)$. Notar que si hacemos insersiones por segmentos la complejidad es $O(\log^2 N)$.


\begin{extended}
    \subsubsection*{Propiedad Transcendente (\textit{Transcending Property})}
    El Li Chao Tree funciona para cualquier familia de funciones continuas $\{f_i\}$ definidas sobre el dominio de interés que cumplan la \textbf{propiedad de intersección única} o \textit{transcending property}:
    \begin{quote}
    \textit{Para cualquier par de funciones distintas $f_i$ y $f_j$, la función diferencia $d(x) = f_i(x) - f_j(x)$ cambia de signo a lo sumo una vez en el intervalo de evaluación $[x_{\min}, x_{\max}]$.}
    \end{quote}
    Esto garantiza que, al comparar dos funciones en el extremo izquierdo \texttt{lv} y en el punto medio \texttt{mv}, la decisión sobre en qué mitad del intervalo se encuentra la intersección permanece binaria, permitiendo descartar una subrama en $\mathcal{O}(1)$ por nivel.

    \subsubsection*{Mecánica de Inserción en \texttt{addFunction}}
    Para un nodo $v$ que cubre el rango de índices $[l, r)$ con punto medio $m = \lfloor(l+r)/2\rfloor$:
    \begin{enumerate}
        \item Se evalúan la nueva función $f$ y la función almacenada $g = \texttt{treefunc}[v]$ en el punto medio $x_m = A[m]$.
        \item La función con mayor valor en $x_m$ se conserva en el nodo $v$, mientras que la función perdedora se reasigna a $f$.
        \item Se evalúa cuál función es superior en el extremo izquierdo $x_l = A[l]$: si el ganador en $x_l$ es distinto del ganador en $x_m$, la intersección se ubica en $[l, m)$ y $f$ desciende al hijo izquierdo; de lo contrario, desciende al hijo derecho $[m, r)$.
    \end{enumerate}
    \textit{Nota:} Para consultas de mínimo ($\min_i f_i(a)$), invierta los operadores relacionales (\texttt{>}) por (\texttt{<}) y reemplace el valor neutro nulo por $+\infty$.
\end{extended}



\algoinfo{ \log N }{ N }{  }
\inputcode{cpp}{latex_src/dp/li_chao_tree.cpp}{li_chao_tree.cpp}